\chapter{Аналитическая часть}
В данной части формализованы объекты сцены, а также рассмотрены основные методы решения поставленных задач.

\section{Городская среда}

Структура генерации городской среды основана на двух типах дорог: внешних, образующих границы кварталов, и внутренних, проходящих внутри кварталов. Вдоль внутренних дорог производится автоматическая расстановка зданий с учетом заданных параметров размера города и числа начального числа псевдослучайной последовательности.

Для визуализации городской среды используются следующие объемные модели.
\begin{itemize}
    \item Дорога -- связанный с другими набор полигонов, расположенных в горизонтальной плоскости. Проспекты -- задают границы кварталов, внутриквартальные дороги образуют случайный набор дорог внутри квартала.
    \item Дом представляет собой набор полигонов, содержащий набор стен и крыши. Отдельно описан первый этаж, содержащий вход, типовой этаж, который одинаково описывает 2..N этаже, где N -- количество этажей дома и крышу. Располагаются вдоль внутриквартальных дорог.
    \item Деревья -- располагаемые вдоль внутриквартальных дорог объемные объекты, с горизонтальным сечением в виде правильного восьмиугольника.
\end{itemize}
\section{Описание алгоритмов}
\subsection{Методы представления трехмерных объектов}
Существует несколько основных подходов к представлению трехмерных объектов в компьютерной графике. Первый метод -- сплошное представление, при котором объект описывается аналитическими уравнениями или неявными функциями. Данный подход обеспечивает высокую точность геометрического описания, но требует значительных вычислительных ресурсов для рендеринга.

Второй метод -- представление на основе вокселей, где пространство разбивается на равномерную трехмерную сетку, и каждый элемент хранит информацию о принадлежности к объекту. Этот подход эффективен для объемных данных, но имеет высокую память для хранения информации.
Третий метод -- каркасное представление, в котором объект описывается набором вершин и соединяющих их рёбер. Данный метод позволяет быстро визуализировать геометрию, но не обеспечивает реалистичного отображения поверхностей.

Четвертый метод -- полигональное представление, при котором поверхность объекта аппроксимируется набором плоских многоугольников. По соображениям вычислительной эффективности и простоты обработки в компьютерной графике преимущественно используются треугольные полигоны, поскольку любой треугольник всегда является плоским и легко поддается обработке. Математически полигональная модель объекта может быть представлена как множество вершин $V = \{v_1, v_2, \dots, v_n\}$ и множество граней $F = \{f_1, f_2, \dots, f_m\}$, где каждая грань $f_i$ определяется набором индексов вершин~\cite{rogers_mathematical}.
Полигональное представление также можно дополнить при помощи указания нормалей поверхности, это позволяет отсекать поверхности выпуклых объемных моделей, которые повернуты от направления наблюдателя, и, соответственно не могут быть видны с текущего ракурса.

В рамках визуализации городской среды полигональное представление с треугольными полигонами является оптимальным выбором, поскольку позволяет эффективно моделировать сложные архитектурные формы зданий и инфраструктуры при приемлемых вычислительных затратах, также будет использовано предвычисление нормалей полигонов для обеспечения отсечений.

\subsection{Методы проецирования трехмерных объектов на двумерную плоскость}
Для отображения трехмерных объектов на двумерном экране используются различные методы проецирования. Ортогональная проекция сохраняет параллельность линий и пропорции объектов, но не создает эффекта глубины. Математически ортогональная проекция на плоскость $xy$ описывается матрицей:

\begin{equation}
P_{ortho} = \begin{bmatrix}
1 & 0 & 0 & 0 \\
0 & 1 & 0 & 0 \\
0 & 0 & 0 & 0 \\
0 & 0 & 0 & 1
\end{bmatrix}.
\end{equation}
Перспективная проекция создает более реалистичное изображение, имитируя работу человеческого глаза, где удаленные объекты кажутся меньше. Данная проекция описывается нелинейным преобразованием, которое в однородных координатах может быть представлено матрицей:

\begin{equation}
P_{persp} = \begin{bmatrix}
\frac{2 \cdot n}{r - l} & 0 & \frac{r + l}{r - l} & 0 \\
0 & \frac{2 \cdot n}{t - b} & \frac{t + b}{t - b} & 0 \\
0 & 0 & -\frac{f + n}{f - n} & -\frac{2 \cdot f \cdot n}{f - n} \\
0 & 0 & -1 & 0
\end{bmatrix},
\end{equation}

где $l, r, b, t$ -- границы усеченной пирамиды видимости по осям $x$ и $y$, $n$ и $f$ -- расстояния до ближней и дальней плоскостей отсечения соответственно.

Для корректного отображения взаимного расположения объектов необходимо решить проблему видимости. Наиболее распространенные методы решения этой задачи -- Z-буфер и алгоритм художника. Z-буфер использует дополнительный буфер, хранящий глубину для каждого пикселя изображения. Для каждого фрагмента проверяется его глубина $z_{frag}$ по отношению к текущему значению в буфере $z_{buffer}$:


\begin{equation}
\text{if } z_{frag} < z_{buffer}[x,y] \text{ then }
\begin{cases}
z_{buffer}[x,y] = z_{frag} \\
\text{color}[x,y] = \text{fragment\_color}
\end{cases}.
\end{equation}
Алгоритм художника представляет собой метод решения проблемы видимости, основанный на аналогии с техникой рисования художника, который сначала изображает дальний фон, затем средние планы и в конце -- ближние объекты. Суть алгоритма заключается в сортировке всех полигонов сцены по убыванию расстояния от точки наблюдения и последующей отрисовке в порядке от наиболее удаленных к ближайшим. Математически это выражается как упорядочение множества полигонов $P = \{p_1, p_2, \dots, p_m\}$ по ключу $d_i = \|c - p_i.center\|$, где $c$ -- позиция камеры, а $p_i.center$ -- центр полигона $p_i$. Однако данный алгоритм имеет принципиальные ограничения: он не может корректно обрабатывать случаи взаимного пересечения полигонов или циклических перекрытий, например, когда полигон A перекрывает B, B перекрывает C, а C перекрывает A. Для преодоления этих ограничений применяются дополнительные техники, такие как разбиение полигонов на более мелкие части или использование гибридных подходов.

Альтернативные методы решения проблемы видимости включают алгоритм Варнока и алгоритм Робертса. Алгоритм Варнока основан на рекурсивном разбиении экрана на подобласти до тех пор, пока в каждой области не останется тривиальная конфигурация для определения видимости. Алгоритм работает в объектном пространстве и на каждом шаге рекурсии анализирует текущую область изображения: если в области находится только один полигон, он отрисовывается полностью; если полигонов несколько и они не перекрываются, выбирается ближайший; если конфигурация сложная, область разделяется на четыре подобласти, и процесс повторяется. Данный подход эффективен для сцен с большим количеством однородных областей, но имеет высокую вычислительную сложность для детализированных изображений.

Алгоритм Робертса является одним из первых методов удаления невидимых линий и поверхностей. Он работает на уровне отдельных рёбер каркасной модели и основан на аналитической геометрии. Алгоритм последовательно проверяет каждое ребро объекта на видимость, анализируя его взаимное расположение с другими объектами сцены. Для каждого ребра строится уравнение плоскости, содержащей это ребро, и проверяется, не перекрывается ли оно другими поверхностями. Основное преимущество алгоритма -- точность определения видимых контуров, однако его практическое применение ограничено высокой вычислительной сложностью $O(n^2)$ для $n$ объектов и сложностью обработки неконвексных моделей.

Для визуализации городской среды метод Z-буфера является предпочтительным, так как обеспечивает корректное отображение взаимного расположения зданий и инфраструктуры без необходимости сложной сортировки геометрии~\cite{rogers_computer_graphics}.
\subsection{Методы освещения и расчета теней}
Для создания реалистичного изображения необходимо моделировать взаимодействие света с поверхностями объектов. В компьютерной графике существует несколько основных подходов к расчету освещения.

Существует несколько классических моделей освещения, отличающихся сложностью вычислений и визуальным качеством. Модель Фонга предполагает интерполяцию нормалей к поверхности в пределах полигона с последующим расчетом освещения для каждого пикселя. Интенсивность освещения в точке рассчитывается по формуле:

\begin{equation}
I = I_a \cdot k_a + I_l \cdot k_d \cdot (\mathbf{N} \cdot \mathbf{L}) + I_l \cdot k_s \cdot (\mathbf{R} \cdot \mathbf{V})^n,
\end{equation}

где $I_a$, $I_l$ -- интенсивности фонового и направленного освещения соответственно, $k_a$, $k_d$, $k_s$ -- коэффициенты фонового, диффузного и зеркального отражения, $\mathbf{N}$ -- нормаль к поверхности, $\mathbf{L}$ -- направление к источнику света, $\mathbf{R}$ -- отраженный вектор, $\mathbf{V}$ -- направление к наблюдателю, $n$ -- степень зеркального отражения~\cite{phong_model}.

Закраска по Гуро использует интерполяцию цвета между вершинами полигона. Освещение рассчитывается только в вершинах, а затем цвет интерполируется по поверхности полигона. Этот метод менее ресурсоемкий, чем модель Фонга, но может приводить к артефактам в виде ``световых пятен`` на гладких поверхностях~\cite{gouraud_shading}.

Рей-трейсинг представляет собой метод глобального освещения, моделирующий физическое поведение света. Алгоритм трассирует лучи от наблюдателя через каждый пиксель изображения, рассчитывая отражения, преломления и тени. Математически интенсивность пикселя определяется рекурсивным соотношением:

\begin{equation}
I_p = I_{direct} + \sum_{i=1}^{n} k_{r_i} \cdot I_{reflected_i} + \sum_{j=1}^{m} k_{t_j} \cdot I_{refracted_j},
\end{equation}

где $I_{direct}$ -- прямое освещение, $k_{r_i}$, $k_{t_j}$ -- коэффициенты отражения и преломления, $I_{reflected_i}$, $I_{refracted_j}$ -- интенсивности отраженных и преломленных лучей~\cite{whitted_raytracing}.

Для визуализации городской среды методы с интерполяцией освещения не были выбраны, поскольку здания преимущественно имеют явные углы и геометрические особенности, где интерполяция нормалей или цвета не требуется. Рей-трейсинг, несмотря на высокое качество результата, требует значительных вычислительных ресурсов и не подходит для интерактивной визуализации больших городских моделей~\cite{realtime_rendering}.

Был выбран метод однотонной закраски с расчетом диффузной яркости, поскольку он обеспечивает хороший баланс между вычислительной сложностью и визуальным качеством. Диффузная составляющая зависит от угла падения света на поверхность и описывается законом Ламберта:

\begin{equation}
I_{diffuse} = I_l \cdot k_d \cdot \max(0, \mathbf{N} \cdot \mathbf{L}).
\end{equation}
Для расчета теней используется метод карты теней, который включает два основных шага.
\begin{enumerate}
\item Рендеринг сцены с точки зрения источника света с сохранением буфера глубины $z_{light}$.
\item Рендеринг сцены с точки зрения наблюдателя с проверкой видимости каждого фрагмента из точки зрения источника света.
\end{enumerate}
% #todo проверить гостовский вариант написания гдешек
\begin{equation}
\text{if } z_{frag\_to\_light} > z_{light}[u,v] + \epsilon \text{ then fragment in shadow},
\end{equation}

где $u,v$ -- координаты в текстурном пространстве карты теней, $\epsilon$ -- небольшая константа для борьбы с артефактами~\cite{shadow_mapping_technique}.

\subsection{Выбранные методы для визуализации городской среды}
Для решения задачи трехмерной визуализации городской среды были выбраны следующие методы:

\textbf{Полигональное представление с треугольными полигонами} -- данное представление обеспечивает оптимальное соотношение между точностью геометрического описания и вычислительной эффективностью. Здания и инфраструктура города могут быть эффективно аппроксимированы треугольными сетками, что позволяет использовать современные графические API для ускорения рендеринга.

\textbf{Z-буфер для определения видимости} -- метод Z-буфера обеспечивает корректное отображение взаимного расположения объектов без необходимости предварительной сортировки геометрии. Алгоритм работает в два этапа: сначала вычисляется глубина каждого фрагмента, затем производится сравнение с текущим значением в буфере глубины для принятия решения о видимости фрагмента.

\textbf{Однотонная закраска с диффузной составляющей} -- для создания реалистичного освещения используется модель, учитывающая диффузное отражение света от поверхностей зданий. Интенсивность освещения рассчитывается по формуле:
% #todo тут бы тоже обосновать
\begin{equation}
I = (\mathbf{N} \cdot \mathbf{L}) \cdot I_{light},
\end{equation}

где $\mathbf{N}$ -- нормаль к поверхности, $\mathbf{L}$ -- направление к источнику света, $I_{light}$ -- интенсивность источника света.

\textbf{Метод карты теней} -- для расчета теней от зданий и других объектов используется алгоритм карты теней. Процесс включает генерацию карты глубины с точки зрения источника света и последующую проверку нахождения фрагментов в тени при рендеринге сцены с точки зрения наблюдателя. Метод обеспечивает реалистичное отображение теней, отбрасываемых зданиями на дороги и другие объекты городской среды. Кроме того, карта теней использует в себе Z-буффер, что увеличивает повторное использование кода.


\section{Аффинные преобразования для трехмерных моделей}
Аффинные преобразования являются фундаментальными операциями в компьютерной графике, позволяющими выполнять геометрические преобразования моделей с сохранением коллинеарности точек и отношения расстояний вдоль прямых линий . В общем виде аффинное преобразование можно представить как комбинацию линейного преобразования и переноса:
\begin{equation}
	\mathbf{P}' = \mathbf{M} \cdot \mathbf{P} + \mathbf{T},
\end{equation}

где $\mathbf{P}$ -- исходная точка, $\mathbf{P}'$ -- преобразованная точка, $\mathbf{M}$ -- матрица линейного преобразования размером $3 \times 3$, $\mathbf{T}$ -- вектор переноса.

Для удобства вычислений в компьютерной графике используется однородная система координат, что позволяет представить любое аффинное преобразование в виде единой матрицы $4 \times 4$:
\begin{equation}
	\mathbf{A} = \begin{bmatrix}
		a_{11} & a_{12} & a_{13} & t_x \\
		a_{21} & a_{22} & a_{23} & t_y \\
		a_{31} & a_{32} & a_{33} & t_z \\
		0 & 0 & 0 & 1
	\end{bmatrix},
\end{equation}

где верхняя левая подматрица $3 \times 3$ описывает линейное преобразование, а правый столбец задает вектор переноса.

Основными типами аффинных преобразований, применяемых к трехмерным моделям, являются:

\subsection{Перенос}
Преобразование переноса смещает все точки модели на заданный вектор $\mathbf{T} = (t_x, t_y, t_z)$:
\begin{equation}
	\mathbf{T}(t_x, t_y, t_z) = \begin{bmatrix}
		1 & 0 & 0 & t_x \\
		0 & 1 & 0 & t_y \\
		0 & 0 & 1 & t_z \\
		0 & 0 & 0 & 1
	\end{bmatrix},
\end{equation}

где $t_x$, $t_y$, $t_z$ -- компоненты вектора переноса вдоль соответствующих осей координат.

\subsection{Масштабирование}
Преобразование масштабирования изменяет размеры модели относительно начала координат:
\begin{equation}
	\mathbf{S}(s_x, s_y, s_z) = \begin{bmatrix}
		s_x & 0 & 0 & 0 \\
		0 & s_y & 0 & 0 \\
		0 & 0 & s_z & 0 \\
		0 & 0 & 0 & 1
	\end{bmatrix},
\end{equation}

где $s_x$, $s_y$, $s_z$ -- коэффициенты масштабирования по осям $x$, $y$ и $z$ соответственно.

\subsection{Вращение}
Вращение модели вокруг оси $x$ на угол $\theta$ описывается матрицей:
\begin{equation}
	\mathbf{R}_x(\theta) = \begin{bmatrix}
		1 & 0 & 0 & 0 \\
		0 & \cos \theta & -\sin \theta & 0 \\
		0 & \sin \theta & \cos \theta & 0 \\
		0 & 0 & 0 & 1
	\end{bmatrix}.
\end{equation}
Вращение вокруг оси $y$:
\begin{equation}
	\mathbf{R}_y(\theta) = \begin{bmatrix}
		\cos \theta & 0 & \sin \theta & 0 \\
		0 & 1 & 0 & 0 \\
		-\sin \theta & 0 & \cos \theta & 0 \\
		0 & 0 & 0 & 1
	\end{bmatrix}.
\end{equation}
Вращение вокруг оси $z$:
\begin{equation}
	\mathbf{R}_z(\theta) = \begin{bmatrix}
		\cos \theta & -\sin \theta & 0 & 0 \\
		\sin \theta & \cos \theta & 0 & 0 \\
		0 & 0 & 1 & 0 \\
		0 & 0 & 0 & 1
	\end{bmatrix},
\end{equation}

где $\theta$ -- угол поворота в радианах.

Сложные преобразования формируются путём композиции базовых аффинных преобразований. Поскольку матричное умножение не коммутативно, порядок применения преобразований существенно влияет на результат. Для преобразования точки $\mathbf{P}$ с помощью последовательности преобразований $\mathbf{A}_1$, $\mathbf{A}_2$, $\ldots$, $\mathbf{A}_n$ используется формула:
\begin{equation}
	\mathbf{P}' = \mathbf{A}_n \cdot \mathbf{A}_{n-1} \cdot \ldots \cdot \mathbf{A}_1 \cdot \mathbf{P},
\end{equation}

где преобразования применяются в порядке справа налево.

Аффинные преобразования обладают важным свойством инвариантности относительно выпуклых комбинаций точек. Если точка $\mathbf{Q}$ является выпуклой комбинацией точек $\mathbf{P}_1$ и $\mathbf{P}_2$:
\begin{equation}
	\mathbf{Q} = \alpha \cdot \mathbf{P}_1 + (1 - \alpha) \cdot \mathbf{P}_2, \quad 0 \leq \alpha \leq 1
\end{equation}
то после применения аффинного преобразования $\mathbf{A}$ сохраняется соотношение:
\begin{equation}
	\mathbf{A}(\mathbf{Q}) = \alpha \cdot \mathbf{A}(\mathbf{P}_1) + (1 - \alpha) \cdot \mathbf{A}(\mathbf{P}_2).
\end{equation}
Это свойство гарантирует сохранение формы и пропорций модели при преобразованиях~\cite{kaji}.

\section*{Вывод}
В данной главе были рассмотрены основные методы представления трехмерных объектов, проецирования их на двумерную плоскость и расчета освещения, математические основы аффинных преобразований для трехмерных моделей, включая перенос, масштабирование и вращение в однородных координатах. Для решения задачи визуализации городской среды были выбраны: полигональное представление с треугольными полигонами, метод Z-буфера для определения видимости, однотонная закраска с расчетом диффузной яркости и метод карты теней для расчета теней. Данный набор методов обеспечивает оптимальное соотношение между вычислительной сложностью и визуальным качеством результата.